\worksheettitle
  {Домашняя практика}
  {\modulemark{M05-H} Ломаная и многоугольник}

\studentnameblock

\begin{practicebox}[1. Ломаная]
  Для ломаной $A-B-C-D-E$:
  \begin{enumerate}[label=\alph*)]
    \item число вершин \answerblank[15mm];
    \item число звеньев \answerblank[15mm];
    \item звенья \answerblank[85mm].
  \end{enumerate}
\end{practicebox}

\begin{practicebox}[2. Замкните границу]
  Какое звено нужно добавить к ломаной $K-L-M-N$, чтобы получить замкнутую
  границу? \answerblank[25mm]

  Как будет называться многоугольник по числу сторон? \answerblank[38mm]
\end{practicebox}

\begin{practicebox}[3. Элементы многоугольника]
  Для пятиугольника $PQRST$ запишите:
  \begin{enumerate}[label=\alph*)]
    \item все вершины: \answerblank[80mm];
    \item все стороны: \answerblank[90mm];
    \item ещё одно допустимое имя при последовательном обходе:
      \answerblank[35mm].
  \end{enumerate}
\end{practicebox}

\begin{warningbox}[4. Исправьте ошибку]
  «У ломаной из шести вершин всегда шесть звеньев».
  Назовите причину ошибки и уточните утверждение отдельно для незамкнутой
  и замкнутой ломаной. \writinglines{3}
\end{warningbox}
